%-------------------------
% JPL Scientist IV cover letter for Nhat-Minh Nguyen
% Adapted from an entry-level LaTeX cover-letter template.
%------------------------

\documentclass[11pt,letterpaper]{article}

\usepackage{latexsym}
\usepackage[empty]{fullpage}
\usepackage{titlesec}
\usepackage{marvosym}
\usepackage[usenames,dvipsnames]{color}
\usepackage{verbatim}
\usepackage[hidelinks]{hyperref}
\usepackage{fancyhdr}
\usepackage{multicol}
\usepackage{csquotes}
\usepackage{tabularx}
\hypersetup{colorlinks=true,urlcolor=black}
\usepackage[11pt]{moresize}
\usepackage{setspace}
\usepackage{fontspec}
\usepackage[inline]{enumitem}
\usepackage{ragged2e}
\usepackage{anyfontsize}

\usepackage[letterpaper,margin=0.68in]{geometry}
\setlength{\footskip}{8pt}

\setmainfont[
BoldFont=SourceSansPro-Semibold.otf,
ItalicFont=SourceSansPro-RegularIt.otf
]{SourceSansPro-Regular.otf}

\pagestyle{fancy}
\fancyhf{}
\fancyfoot{}
\renewcommand{\headrulewidth}{0pt}
\renewcommand{\footrulewidth}{0pt}

\urlstyle{same}
\raggedbottom
\raggedright
\setlength{\tabcolsep}{0in}

\definecolor{UI_blue}{RGB}{32, 64, 151}
\definecolor{HF_color}{RGB}{179, 179, 179}

\titleformat{\section}{
\color{UI_blue} \scshape \raggedright \large
}{}{0em}{}[\vspace{-0.7cm} \hrulefill \vspace{-0.2cm}]

\begin{document}

\begin{center}
    {\Huge Nhat-Minh Nguyen} \\
    \vspace{0.06cm}
    {\color{UI_blue} \Large Scientist IV, Roman and Euclid Cosmology} \\
    \vspace{0.08cm}
    \large
    \href{mailto:nhat.minh.nguyen@ipmu.jp}{nhat.minh.nguyen@ipmu.jp}
    \quad | \quad
    \href{https://linkedin.com/in/minhmpa}{linkedin.com/in/minhmpa}
    \quad | \quad
    \href{https://github.com/MinhMPA}{github.com/MinhMPA}
    \quad | \quad
    \href{https://orcid.org/0000-0002-2542-7233}{ORCID}

\vspace{0.05cm}
{\color{UI_blue} \hrulefill}
\end{center}

\justify
\setlength{\parindent}{0pt}
\setlength{\parskip}{8.5pt}
\vspace{-0.05cm}

\today

Search Committee \\
Jet Propulsion Laboratory

Dear Search Committee,

I am applying for the Scientist IV position in Roman and Euclid cosmology because it is the kind of problem my research has been moving toward: how to turn precision maps of the late Universe into reliable tests of the dark sector. I work at the interface of nonlinear structure formation, statistical inference, and galaxy surveys. My work develops forward-modeling and field-level Bayesian methods for large-scale structure, with a simple goal: to keep as much information as possible from cosmic structure while keeping the inference physically controlled.

Roman and Euclid will push weak lensing and galaxy clustering into a regime where modeling choices matter as much as statistical power. Euclid brings wide-area coverage. Roman brings depth and image quality. Their overlap can become a testing ground for cross-calibration, validation, and joint inference, especially for galaxy bias, intrinsic alignments, photometric redshifts, selection effects, and cross-probe covariance. These are not separate bookkeeping problems. They decide whether precision measurements become convincing statements about dark energy, dark matter, and gravity.

Two parts of my record are especially relevant. First, I have led the team that found evidence for a late-time suppression in the growth of large-scale structure, which strengthened my interest in growth measurements as tests of dark energy and modified gravity. Second, I have pioneered field-level Bayesian inference methods for galaxy surveys, work recognized by the 2024 Buchalter Cosmology Prize. The technical details differ, but the motivation is the same---the cosmic web contains more information than standard summaries usually keep, and the hard part is extracting it without fooling ourselves.

At JPL, I would focus on making these methods useful in actual Roman and Euclid analyses. That means helping shape joint-probe strategy, testing likelihoods with simulations, using OpenUniverse-style mocks where appropriate, building cross-survey consistency tests, and connecting lensing and clustering measurements to dark-sector interpretation. I do not think every Roman or Euclid analysis should immediately become field-level. The value of field-level inference is that it gives us a controlled way to ask what information is being lost, where summary statistics are safe, and which modeling assumptions will fail first as the error bars shrink. I would also bring experience organizing scientific programs, mentoring junior researchers, working across theory and observational communities, and developing PI-led research through competitive grants.

JPL is a natural place for this work because Roman and Euclid require nonlinear structure modeling, statistical inference, and collaborative mission science to meet in the same analysis. I would approach this role with focus and care. These missions are where field-level and forward-modeling methods can stop being only methodological advances and become part of the machinery that makes dark-sector constraints trustworthy. I would welcome the opportunity to discuss how I could contribute as both a scientific leader and a collaborative member of JPL cosmology program.

\vspace{0.18cm}
Sincerely, \\
Nhat-Minh Nguyen

\end{document}
